\section{Background and Motivation}
\label{sec:background}

\subsection{Pangenome Graphs as a Storage Problem}
\label{subsec:bg-format}

GFA is a line-oriented text format~\cite{gfaspec}. Each line is a record of one
type: \textbf{S} carries a segment identifier and its DNA sequence; \textbf{L},
\textbf{J}, and \textbf{C} carry edges; \textbf{H} carries file-level metadata;
and optional \texttt{TAG:TYPE:VALUE} fields extend any record. Traversals are
\textbf{P} (a named path) and \textbf{W} (a haplotype walk with sample,
haplotype, and coordinate metadata). A traversal is an ordered sequence of
oriented node IDs, where the sign records which DNA strand the segment is read
on, and it is the concatenation of those segments' sequences that reconstructs an
individual's chromosome.

From a storage standpoint the format has three properties that matter. It is
row-major: fields of the same type are scattered across lines rather than
collected into columns. It is textual: a node ID that fits in four bytes is
written as up to nine ASCII characters plus a separator. And, most importantly,
it is \emph{explicitly redundant across records}, because two haplotypes that
agree over a megabase still write out that megabase of node IDs twice.

The third property is the one that drives growth. Adding a haplotype to a
pangenome adds sublinearly to topology, since shared sequence merges into existing
nodes, but adds a full new traversal record. Traversals therefore come to dominate
the file: over $90\%$ in general, and $98.7\%$ of a $14$\,GB chromosome~19 graph
at $1{,}000$ haplotypes~\cite{heringer2025human}. It is also the property that
makes the file so compressible, since the redundancy is long-range and highly
structured. General-purpose compressors capture little of it, because their match
windows are megabytes while the repeats span whole chromosome-length walks: gzip
leaves a $9{,}535$\,GiB $1000$-genomes graph at $2{,}231$\,GiB~\cite{siren2022gbz}.
Domain-specific compressors capture much more, reaching $18$--$84\times$ by
modeling traversal structure directly~\cite{siren2022gbz,gfaz2026}.

\subsection{The Materialize-to-Analyze Tax}
\label{subsec:bg-tax}

Compression has therefore already succeeded at the storage layer, and the
bottleneck has moved. What a practitioner does with a pangenome is derive things
from it, repeatedly: a reference-relative variant table to feed the standard
genomics toolchain, a growth curve to decide whether the cohort is saturated, a
presence/absence matrix for an association study, a coverage profile, a pairwise
similarity matrix. Each derivation today runs a separate single-purpose tool, and
each such tool begins by building a representation larger than the answer it
computes.

Table~\ref{tab:motivation} makes the size of that tax concrete on the largest
public human pangenome. The entire biological output of a growth analysis is a
curve of a few thousand numbers. Producing it with the standard tool costs four
hours and $336$\,GB of resident memory, because the tool must first parse
$369$\,GB of text into item tables. The compressed container for the same graph is
$4.5$\,GB and sits on disk untouched for the whole run.

\begin{table}[t]
\centering
\small
\caption{The cost of one analysis on the HPRC v2.1 pangenome graph, at 16 threads.
The answer is a few thousand integers; the established tool reaches it by
rebuilding a representation $75\times$ larger than the compressed form of the same
data. Setup in Section~\ref{subsec:eval-setup}.}
\label{tab:motivation}
\begin{tabular}{@{}lr@{}}
\toprule
Uncompressed GFA text                  & $369$\,GB \\
Compressed container~\cite{gfaz2026}   & $4.5$\,GB \\
\midrule
Panacus \cmd{growth}: time / peak RSS  & $251$\,min / $336$\,GB \\
\sys \cmd{growth}: time / peak RSS     & $41.2$\,s / $13.7$\,GB \\
\midrule
Output size (growth curve)             & $<$$100$\,KB \\
\bottomrule
\end{tabular}
\end{table}

We call this the \emph{materialize-to-analyze tax}, and it has three components,
which matter because they respond to different fixes. There is an \textbf{I/O}
component: the uncompressed bytes must be read from storage. There is a
\textbf{parse and build} component: text must be converted into an in-memory
structure, which is CPU-bound and usually single-pass but rarely parallel. And
there is a \textbf{residency} component: the built structure must stay in memory
for the duration, which is what sets the machine you need and what makes the two
largest graphs simply infeasible for two of the three baselines.

\subsection{Why the Obvious Fixes Are Not Enough}
\label{subsec:bg-alternatives}

Three responses suggest themselves before ours, and it is worth saying why each
falls short, because reviewers of a storage paper will reach for them first.

\paragraph{Decompress to fast local storage.}
The cheapest fix is to keep the archive compressed, expand it to NVMe on demand,
and run the existing tool. This attacks only the I/O component. It leaves parse,
build, and residency untouched, and it adds a full write of the expanded data plus
the capacity to hold it. For the graphs where the problem is acute, residency is
the binding constraint, not bandwidth: \cmd{vg} and \cmd{odgi} fail on the
$358$/$369$\,GB graphs because the built object does not fit, and no storage tier
changes that. We measure this alternative directly in
Section~\ref{subsec:eval-expand} rather than argue about it.

\paragraph{Buy more memory.}
Residency can be bought, up to a point, and the point is closer than it looks.
Baseline peak RSS on these workloads runs $4$--$5\times$ the GFA size for the
graph-resident tools; at the current $369$\,GB that is already beyond a
$503$\,GiB node, and cohort sizes are growing faster than per-node memory. The
approach also does nothing for the analyses' time cost, which is what a scientist
actually waits on when sweeping parameters or reprocessing a new release.

\paragraph{Use a succinct graph index.}
The strongest existing answer is to store the graph in a compressed
\emph{queryable} index rather than as text. GBZ~\cite{siren2022gbz}, built on the
GBWT~\cite{siren2020haplotype,siren2021pangenomics}, does exactly this and is the
closest intellectual relative of our work: it keeps haplotype paths in a
BWT-derived structure that supports navigation without expansion. It is genuinely
good at what it targets, and tools such as \cmd{vg giraffe} depend on it.

The difference is what the structure is optimized for. GBZ is built for
\emph{navigational} queries: locate a pattern, extend a search, enumerate the
haplotypes through a position. Those require an index with rank/select support
over the whole collection, which is global to build, does not admit incremental
extension without rebuilding, and does not preserve full GFA record semantics such
as optional fields. Our workloads are the opposite shape. They are dense sequential
folds that touch every step of every traversal exactly once and need no locate at
all. For that access pattern an ordered stream is the right structure, and the
machinery that makes locate fast is overhead that a fold cannot amortize. The two
designs are complementary rather than competing, and we compare against GBZ-backed
tooling where a common analysis exists (Section~\ref{subsec:eval-gbz}).

\subsection{What Would Have to Be True}
\label{subsec:bg-requirements}

If a query is going to run on the compressed bytes, then the format has to hold up
under a set of demands it was never designed for: a decoder that emits elements
without expanding the whole record, a way to touch one field of the data without
paying for the others, and enough independence between records that many threads
can decode at once. Section~\ref{sec:principles} states these as five properties
and checks them against formats in common use; Section~\ref{sec:container}
describes the container we build on in exactly those terms.
